%=============================================================================!
% 16.F  SOLUTIONS TO THE EXERCISES                                            !
%=============================================================================!
\unnumbered{Chapter~\ref{ch:observables}: Observables}
\label{sec:ob-solutions}

\solutionto{ex:ob-shell}The exact volume is $\frac43\pi[(r + \delta r)^3 -
r^3] = 4\pi r^2\delta r(1 + \delta r/r + \delta r^2/3r^2)$, so the first
term falls short by the fraction $1 - 1/(1 + x) = x/(1 + x)$, with $x =
\delta r/r + \delta r^2/3r^2$, which is close to $x$, and so to $\delta
r/r$, when $\delta r \ll r$. Numerically, at
\SI{1}{\angstrom} it is 0.0484 and at \SI{5}{\angstrom} 0.0099, against
$\delta r/r = 0.05$ and 0.01.

\solutionto{ex:ob-far}$1 - 1/256 = 0.99609$ and $1 - 1/2048 = 0.99951$.
Dividing by $N(N - 1)/V$ instead of $N^2/V$ makes it approach 1. The
coordination number counts partners, $\rho g(r)$ times the volume of each
shell; the new divisor scales $g(r)$ up by $N/(N - 1)$ and goes with the
partners' density $(N - 1)/V$, smaller by the same factor, so the count
is unchanged. Used with $\rho = N/V$ it would add $12/(N - 1) = 0.05$ to
the 12 of the first shell for 256 atoms. In a liquid in a large enough
cell the far value under ASE's normalisation is $1 - S(0)/N$ instead,
within \num{3e-5} of 1 for 2048 atoms of argon.

\solutionto{ex:ob-fcc}With an atom at a corner as the origin, every atom
of the lattice sits at $(a/2)(n_1, n_2, n_3)$ with whole numbers whose sum
is even: all three even at the corners, two odd and one even at the face
centres. Its squared distance is $(a^2/4)(n_1^2 + n_2^2 + n_3^2)$, and the
sum of squares is then even. The smallest value, 2, needs two of the $n$
equal to $\pm1$ and the third 0: three places for the 0 times four
choices of signs, 12 atoms at $\sqrt2\,a/2 = a/\sqrt2$. The next, 4,
cannot come from $(\pm1, \pm1, \pm1)$, whose sum is odd, so it needs one
of them $\pm2$ and the others 0: 6 atoms at the distance $a$.

\solutionto{ex:ob-meanfield}With $g = 1$ beyond $\sigma$,
\begin{equation*}
   \frac{U}{N} = \frac{\rho}{2}\int_\sigma^\infty4\varepsilon\Bigl[
   \Bigl(\frac\sigma r\Bigr)^{12} - \Bigl(\frac\sigma r\Bigr)^6\Bigr]
   4\pi r^2\,\dd r
   = 8\pi\rho\varepsilon\sigma^3\Bigl(\frac19 - \frac13\Bigr)
   = -\frac{16}{9}\pi\rho\sigma^3\varepsilon ,
\end{equation*}
since $\int_\sigma^\infty\sigma^{12}r^{-10}\,\dd r = \sigma^3/9$ and
$\int_\sigma^\infty\sigma^6r^{-4}\,\dd r = \sigma^3/3$. With $\rho\sigma^3
= 0.8$ it is \SI{-46.20}{\milli\electronvolt}, \SI{9.97}{\milli
\electronvolt} above the value from the real $g(r)$. The difference
measures the liquid's structure: the real $g(r)$ crowds partners into the
first shell, where the potential is most negative. The run's own
\SI{-50.09}{\milli\electronvolt} lies higher than \SI{-56.16}{\milli
\electronvolt} because its switched potential leaves out the attraction
beyond $2\sigma$; with that potential the mean field itself gives
\SI{-40.07}{\milli\electronvolt}, and the structure again about
\SI{10}{\milli\electronvolt}.

\solutionto{ex:ob-sq-limits}$S(\vb{G}) = 1 + (1/N)\sum_i\sum_{j\neq
i}\cos(\vb{G}\cdot\vb{r}_{ij})$. For independent positions spread evenly
over the cell, each cosine averages over the cell to zero for $\vb{G}
\neq \vb{0}$ (\cref{sec:md-ewald}), leaving $S(\vb{G}) = 1$. In a perfect
crystal at a reciprocal lattice vector every $\vb{G}\cdot\vb{r}_{ij}$ is a
multiple of $2\pi$, every cosine is 1, and $S(\vb{G}) = 1 + N(N - 1)/N =
N$.

\solutionto{ex:ob-s0}$S(0) = \rho\kB T\kappa_T$ with $\kappa_T = 0.5
\times160.2 = \SI{80.1}{\angstrom\cubed\per\electronvolt}$, so $S(0) = 0.03
\times0.02585\times80.1 = 0.0621$. Twice as compressible, twice $S(0)$: a
softer liquid lets its density fluctuate more, and scatters long waves
more strongly.

\solutionto{ex:ob-flicker}Each period carries the pair out across the
cutoff and back, two events, so $2\times1000/60 = 33.3$ events per
picosecond. With $r_{\mathrm{on}}$ at the old cutoff the bond, once
formed, breaks only if the pair passes $r_{\mathrm{off}}$, so the band
records none if it is wider than the part of the swing beyond the
cutoff. Wherever the cutoff falls within the vibration, a band wider
than the whole swing, from its near end to its far end, records none,
since no swing can then carry the pair across both distances.

\solutionto{ex:ob-keys}With $N = 6$ the first frame's keys are $0\cdot6 + 1
= 1$, $1\cdot6 + 2 = 8$ and $3\cdot6 + 4 = 22$; the second's are 1, 15,
22 and 29. Merging the sorted lists $(1, 8, 22)$ and $(1, 15, 22, 29)$:
1 is in both; 8 is smaller than 15 and only in the first, broken; 15 is
smaller than 22 and only in the second, formed; 22 is in both; 29 is left
in the second, formed. So 2-3 and 4-5 form and 1-2 breaks. The search
starting at atom 0 finds $\{0, 1, 2\}$, then $\{3, 4\}$, then $\{5\}$ in
the first frame, and $\{0, 1\}$ and $\{2, 3, 4, 5\}$ in the second.

\solutionto{ex:ob-hash}In the prism each atom is joined to the two others
of its triangle and to one atom of the other; in the hexagon each is
joined to its two neighbours round the ring and to the opposite atom.
Every atom has three neighbours, all of the same element. In the first
round every atom starts with the same number, and its neighbours' sum is
three times the same scrambled number, so every atom of both gets the
same new number; the same holds in every later round, and the sums over
the molecules are equal. The molecules differ: the prism has triangles of
bonds, and the hexagon with its cross-links has none. Its atoms alternate
round the ring between two sets of three, and each bond, round the ring
or across it, joins one set to the other, so any closed path of bonds
returns to its start after an even number of steps.

\solutionto{ex:ob-poisson}The relative error is $1/\sqrt n$, so 2\% needs
$1/0.02^2 = 2500$ events. Forty events over $100\times\SI{1000}{\pico
\second} = \num{1e5}$ molecule-picoseconds give $(4.0 \pm
0.6)\times10^{-4}$ per picosecond, the error $\sqrt{40}/10^5$.

\solutionto{ex:ob-walk}After $n = k_{\mathrm h}t$ hops $x = \sum_ks_k$
with $s_k = \pm a$ independent, so $\ens{x^2} = \sum_k\ens{s_k^2} = na^2 =
k_{\mathrm h}a^2t$. Comparing with $2Dt$ for $d = 1$, $D = k_{\mathrm
h}a^2/2$.

\solutionto{ex:ob-correlated}The probability $\frac12(1 + c)$ of repeating
and $\frac12(1 - c)$ of reversing give $\ens{s_ks_{k+1}} = a^2[\frac12(1 +
c) - \frac12(1 - c)] = ca^2$, and each step depends on the previous alone,
so $\ens{s_ks_{k+m}} = c^ma^2$. Then, as for \cref{eq:cv-var},
\begin{equation*}
   \ens{x^2} = \sum_k\sum_l\ens{s_ks_l} = a^2\Bigl[n + 2\sum_{m\ge1}(n -
   m)c^m\Bigr] \to na^2\Bigl[1 + \frac{2c}{1 - c}\Bigr] = na^2\,\frac{1 +
   c}{1 - c}
\end{equation*}
for many steps, so $f = (1 + c)/(1 - c)$. Solving $(1 + c)/(1 - c) =
3.46$ gives $c = (3.46 - 1)/(3.46 + 1) = 0.552$.

\solutionto{ex:ob-water}The motion stops being ballistic near $t =
2Dm/\kB T$. With $\SI{1}{\amu} = \SI{103.64}{\electronvolt\femto\second
\squared\per\angstrom\squared}$, $m = 18.015\times103.64$ in those units,
and $\kB T = \SI{0.02620}{\electronvolt}$, this is about
\SI{75}{\femto\second}.

\solutionto{ex:ob-exp}Writing $\tau$ for $\tau_{\mathrm c}$ in this
solution, with $C_{vv}(s) = C_0\mathrm{e}^{-s/\tau}$,
\begin{align*}
   2\int_0^t(t - s)C_0\mathrm{e}^{-s/\tau}\,\dd s
   &= 2C_0\bigl[t\tau(1 - \mathrm{e}^{-t/\tau}) - \tau^2 + (\tau^2 +
   t\tau)\mathrm{e}^{-t/\tau}\bigr]\\
   &= 2C_0\tau\bigl[t - \tau(1 - \mathrm{e}^{-t/\tau})\bigr] ,
\end{align*}
using $\int_0^ts\mathrm{e}^{-s/\tau}\,\dd s = \tau^2 - (\tau^2 +
t\tau)\mathrm{e}^{-t/\tau}$, by parts. For $t \ll \tau$, $1 -
\mathrm{e}^{-t/\tau} = t/\tau - t^2/2\tau^2 + \cdots$, and the bracket is
$t^2/2\tau$, so the MSD is $C_0t^2$, the ballistic form with $C_0 =
3\kB T/m$. For $t \gg \tau$ the bracket is $t - \tau$, the MSD grows as
$2C_0\tau t$, and $6D = 2C_0\tau$ gives $D = C_0\tau/3$, which is also
$\frac13\int_0^\infty C_{vv}$.

\solutionto{ex:ob-eta}The integral of $\ens{\mathsf{P}_{xy}(0)\mathsf{P}_{xy}
(t)}$ is the variance $\sigma_{xy}^2$ times the integral of the normalised
autocorrelation function, which is $\tau_{\mathrm{int}}$ by its definition
(\cref{sec:cv-correlation}). With $V = 256/0.02035 =
\SI{12580}{\angstrom\cubed}$, $\kB T = \SI{0.01163}{\electronvolt}$,
$\sigma_{xy} = 0.0121/160.2 =
\SI{7.55e-5}{\electronvolt\per\angstrom\cubed}$
and $\tau_{\mathrm{int}} = \SI{180}{\femto\second}$, the product is
$\SI{1.11}{\electronvolt\femto\second\per\angstrom\cubed}$, and with
$\SI{1}{\electronvolt\per\angstrom\cubed} = \SI{1.602e11}{\pascal}$ and
$\SI{1}{\femto\second} = \SI{e-15}{\second}$ it is \SI{1.78e-4}{\pascal
\second}, near the \num{1.853e-4} of \cref{sec:ob-vacf}.

\solutionto{ex:ob-ne}$\sigma_{\mathrm{NE}} = e^2\sum_iD_i/V\kB T$ with the
sum $32\times(5.19 + 5.03)\times10^{-4}$\,\si{\angstrom\squared\per\femto
\second} $= \SI{3.27e-2}{\angstrom\squared\per\femto\second} =
\SI{3.27e-7}{\metre\squared\per\second}$, $V = \SI{2.406e-27}{\metre\cubed}$
and $\kB T = \SI{2.053e-20}{\joule}$, and $e^2 = \SI{2.567e-38}{\coulomb
\squared}$: \SI{170}{\siemens\per\metre}, close to the mean of the runs'
own values, 168, since each run's temperature differs.

\solutionto{ex:ob-trace}A sixth of the slope of the full MSD is a third of
the trace, $(1.16 + 1.02 + 0.00)/3 = \SI{0.73e-4}{\angstrom\squared\per
\femto\second}$. It is neither the fast motion in the plane, $1.1\times
10^{-4}$, nor the motion across it, for which no crossing was observed. A single
number hides this distinction; the one-crossing scale is about $10^{-7}$.

\solutionto{ex:ob-euler}$\lvert\mathrm{e}^{\mathrm{i}\theta}\rvert^2 =
\cos^2\theta + \sin^2\theta = 1$. Multiplying $(\cos A + \mathrm{i}\sin
A)(\cos B + \mathrm{i}\sin B) = \cos A\cos B - \sin A\sin B +
\mathrm{i}(\sin A\cos B + \cos A\sin B)$ and equating its real and
imaginary parts with those of $\cos(A + B) + \mathrm{i}\sin(A + B)$ gives
both formulae.

\solutionto{ex:ob-dft-cos}Writing $N$ for $N_{\mathrm s}$ in this solution,
$\cos(2\pi mn/N) = \frac12[\mathrm{e}^{2\pi\mathrm{i}mn/N} +
\mathrm{e}^{-2\pi\mathrm{i}mn/N}]$, so
\begin{equation*}
   X_k = \frac12\sum_n\mathrm{e}^{2\pi\mathrm{i}(m - k)n/N}
   + \frac12\sum_n\mathrm{e}^{-2\pi\mathrm{i}(m + k)n/N} .
\end{equation*}
By the geometric sum of \cref{sec:ob-fourier} the first is $N/2$ when $k =
m$ and zero otherwise, and the second is $N/2$ when $m + k$ is a multiple
of $N$, that is $k = N - m$, and zero otherwise.

\solutionto{ex:ob-fft-cost}Writing $N$ for $N_{\mathrm s}$, if $W(N/2) =
(N/2)\log_2(N/2)$ then $W(N) = N[\log_2N - 1] + N = N\log_2N$, and $W(1)
= 1\times0 = 0$ starts the chain. For $N = 2^{20}$ the ratio
$N^2/N\log_2N = N/20 = 52\,429$.

\solutionto{ex:ob-fold}The Nyquist frequency of \SI{35}{\femto\second} is
\SI{14.29}{\tera\hertz}, and \SI{18.25}{\tera\hertz} folds to $1/\Dt - \nu =
28.57 - 18.25 = \SI{10.32}{\tera\hertz}$. It appears where it is for strides
below $1/(2\times\SI{18.25}{\tera\hertz}) = \SI{27.4}{\femto\second}$.

\solutionto{ex:ob-width}The friction widens each peak to $\gamma/2\pi =
\SI{0.159}{\tera\hertz}$; the window of \SI{10}{\pico\second} blurs it over
about $1/t_{\mathrm c} = \SI{0.10}{\tera\hertz}$. The friction does more,
though the two are comparable.

\solutionto{ex:ob-vdos-zero}$12mD/\kB T = 12\times15.999\times
\num{103.64}\times\num{5.24e-4}/0.02620 = \SI{398}{\femto\second}$, which
is 0.398 per terahertz.

\solutionto{ex:ob-fd}$F = \frac12kr_{\mathrm h}^2 - \kB T\ln r_{\mathrm h}$ up to a constant, so
$\dd F/\dd r_{\mathrm h} = kr_{\mathrm h} - \kB T/r_{\mathrm h}$, zero at $r_{\mathrm h}^2 = \kB T/k$, and $\dd^2F/\dd r_{\mathrm h}^2 =
k + \kB T/r_{\mathrm h}^2 > 0$, a minimum. At \SI{300}{\kelvin}, $r_{\mathrm h} = \sqrt{0.02585/
0.9785} = \SI{0.163}{\angstrom}$.

\solutionto{ex:ob-crossings}The free energy difference is $-\kB T\ln$ of
the ratio of the times spent in the two wells. In the two-state approximation, each crossing starts a
new, independent stay, so with $n$ crossings each well is visited about
$n/2$ times. The stays have a spread about as large as their mean, as
the waits between Poisson events do (\cref{sec:ob-reactions}), so each
well's total time is known to about $1/\sqrt{n/2} = \sqrt{2/n}$ of
itself, and their ratio, the two errors combined, to $\sqrt2\sqrt{2/n} =
2/\sqrt n$. A relative change $\delta$ in the ratio changes its logarithm
by $\delta$, so the difference is known to about $2\kB T/\sqrt n$, and
$0.1\,\kB T$ needs 400 crossings.
